\begin{proof}[Proof of \cref{obj:thm:unrestricted-uniform-mixed-count-certificate}]
We take \(\overline C=10^{30}\). We first compare the deterministic envelope with the rate, and then establish the risk bound.

\begin{enumerate}
\item Define the untruncated rate expression by
\[
r^\circ_{n,d,q}
:=N^{-1}+\left(\frac{d}{N\ell}\right)^2,
\qquad
N=nq,\quad m=\frac n6,\quad N_0=\frac N6,
\quad \ell=\log(e+N).
\]
The domain assumptions give
\[
0<N\le n,\qquad m>0,\qquad N_0>0,
\qquad r_{n,d,q}=\min\{1,r^\circ_{n,d,q}\}.
\]
If \(N\le1\), then \(N^{-1}\ge1\), so
\[
\mathfrak U_{n,d,q}=r_{n,d,q}=1
\le10^{30}r_{n,d,q}.
\]
For the envelope comparison, henceforth suppose \(N>1\).
% lean: CausalSmith.Stat.MarRareqLogfrontier.riskEnvelope_le_frontierRate

\item The logarithmic scale and stream sizes satisfy
\[
1<\ell,\qquad e^\ell=e+N,\qquad e^{-\ell}\le N^{-1},
\qquad
\frac1{N_0}=6N^{-1},\qquad
\frac1m=\frac6n\le6N^{-1}.
\]
Indeed, \(e+N>e\), and inversion of \(e^\ell=e+N>N\) gives the exponential bound.

The elementary numerical bound for \(\log2\) gives
\[
\gamma_0=\log2-\frac12\ge\frac18.
\]
For every real \(t_{\mathrm{dec}}>0\), the inequality
\(t_{\mathrm{dec}}\le e^{t_{\mathrm{dec}}}\) implies
\(e^{-t_{\mathrm{dec}}}\le t_{\mathrm{dec}}^{-1}\). Applying this with
\(t_{\mathrm{dec}}=n\gamma_0\) yields
\[
e^{-n\gamma_0}
\le\frac1{n\gamma_0}
\le\frac8n
\le8N^{-1}.
\]
% lean: CausalSmith.Stat.MarRareqLogfrontier.poisson_tail_rate

\item Suppose first that \(\ell\ge128\) and \(d\ge\sqrt N\). On this branch,
\[
k=\lfloor\ell/64\rfloor,\qquad B=256\ell,
\qquad
\frac{\ell}{128}\le k\le\frac{\ell}{64}.
\]
For completeness, for every real \(t_{\mathrm{floor}}\ge1\),
\[
\lfloor t_{\mathrm{floor}}\rfloor\ge\frac{t_{\mathrm{floor}}}{2}.
\]
If \(t_{\mathrm{floor}}<2\), its floor is at least one; if
\(t_{\mathrm{floor}}\ge2\), use
\(\lfloor t_{\mathrm{floor}}\rfloor>t_{\mathrm{floor}}-1
\ge t_{\mathrm{floor}}/2\). This proves the lower bound on \(k\).

The inequality \(2t_{\mathrm{exp}}\le e^{t_{\mathrm{exp}}}\), for real
\(t_{\mathrm{exp}}\), gives
\[
\ell\le16e^{\ell/32},
\qquad
\ell^4\le65536e^{\ell/8}.
\]
Since \(e<3\) and \(N>1\),
\[
e^\ell=e+N\le4N\le4d^2,
\qquad e^{\ell/2}\le2d.
\]
Consequently,
\[
\begin{aligned}
(e^{\ell/8}+1)\ell^4
&\le2e^{\ell/8}\,65536e^{\ell/8}\\
&=131072e^{\ell/4}
\le131072e^{\ell/2}
\le262144d,
\end{aligned}
\qquad
\ell\le e^{\ell/2}\le2d.
\]
% lean: CausalSmith.Stat.MarRareqLogfrontier.needle_growth

\item On the same branch, the preceding bounds imply
\[
e^{-16\ell}\le N^{-1},
\qquad e^{-32\ell}\le N^{-1},
\qquad N^{-1}\le\frac{2d}{N\ell}.
\]
Moreover,
\[
\begin{aligned}
\frac{4dB}{N_0k^2}
&=\frac{6144d\ell}{Nk^2}\\
&\le\frac{6144d\ell}{N(\ell/128)^2}
=100663296\,\frac{d}{N\ell}.
\end{aligned}
\]
Thus
\[
\frac{4dB}{N_0k^2}+3e^{-16\ell}
\le100663302\,\frac{d}{N\ell},
\]
and, because both sides are nonnegative,
\[
b_{n,d,q}^2
\le40532401478172816
\left(\frac{d}{N\ell}\right)^2.
\]
% lean: CausalSmith.Stat.MarRareqLogfrontier.riskEnvelope_le_frontierRate

\item Since \(N_0\le m\) and \(B\ge1\),
\[
\frac{B}{mN_0}
\le\frac{B}{N_0^2}
\le\frac{B^2}{N_0^2}.
\]
Therefore
\[
\frac{B^2}{N_0^2}+\frac{B}{mN_0}
\le\frac{2B^2}{N_0^2}
=4718592\,\frac{\ell^2}{N^2}.
\]
The growth estimate above also gives
\[
\frac{d(e^{\ell/8}+1)\ell^2}{N^2}
=
\frac{d(e^{\ell/8}+1)\ell^4}{N^2\ell^2}
\le262144\left(\frac{d}{N\ell}\right)^2.
\]
Combining these inequalities,
\[
64d(e^{\ell/8}+1)
\left(\frac{B^2}{N_0^2}+\frac{B}{mN_0}\right)
\le79164837199872
\left(\frac{d}{N\ell}\right)^2.
\]
The remaining variance terms satisfy
\[
8\left(\frac4{N_0}+\frac1m\right)\le240N^{-1},
\]
and, using \(N^{-1}\le1\) and \(1/m\le6\),
\[
32e^{-32\ell}\left(1+\frac1m\right)\le224N^{-1}.
\]
Together with the prefix-tail bound, this proves
\[
b_{n,d,q}^2+v_{n,d,q}+4e^{-n\gamma_0}
\le40611566315372688
\left(\frac{d}{N\ell}\right)^2+496N^{-1}
\le10^{30}r^\circ_{n,d,q}.
\]
If \(r^\circ_{n,d,q}\le1\), take the minimum with four on the left.
If \(r^\circ_{n,d,q}>1\), use
\(\mathfrak U_{n,d,q}\le4\le10^{30}\). In either case,
\[
\mathfrak U_{n,d,q}\le10^{30}r_{n,d,q}.
\]
% lean: CausalSmith.Stat.MarRareqLogfrontier.riskEnvelope_le_frontierRate

\item On the remaining branch, \(\ell<128\) or \(d<\sqrt N\). Since \(e\ge1\),
\[
0\le\frac{8d}{eN_0}\le\frac{48d}{N},
\qquad
\left(\frac{8d}{eN_0}\right)^2
\le2304\left(\frac dN\right)^2.
\]
If \(\ell<128\), then
\[
\left(\frac dN\right)^2
=\ell^2\left(\frac{d}{N\ell}\right)^2
\le16384\left(\frac{d}{N\ell}\right)^2.
\]
If \(d<\sqrt N\), then \(d^2\le N\), so
\[
\left(\frac dN\right)^2\le N^{-1}.
\]
Hence throughout this branch,
\[
\left(\frac dN\right)^2
\le N^{-1}+16384\left(\frac{d}{N\ell}\right)^2.
\]
Also,
\[
4\left(\frac4{N_0}+\frac1m\right)\le120N^{-1}.
\]
It follows that
\[
\begin{aligned}
\left(\frac{8d}{eN_0}\right)^2
+4\left(\frac4{N_0}+\frac1m\right)+4e^{-n\gamma_0}
&\le2456N^{-1}
+37748736\left(\frac{d}{N\ell}\right)^2\\
&\le10^{30}r^\circ_{n,d,q}.
\end{aligned}
\]
The same split at \(r^\circ_{n,d,q}=1\) proves
\[
\mathfrak U_{n,d,q}\le10^{30}r_{n,d,q}.
\]
This completes the deterministic comparison over the stated domain.
% lean: CausalSmith.Stat.MarRareqLogfrontier.riskEnvelope_le_frontierRate

\item Fix now \(P\in\mathcal M^+_{n,d,q}\). By
\cref{obj:lem:unrestricted-cell-identification,obj:def:ate-functional},
\[
\Phi(P)=\tau(P)
=P(Y(1)=1)-P(Y(0)=1)\in[-1,1].
\]
The auxiliary output in \cref{obj:def:mixed-count-estimator} lies in
\([-1,1]\). Its conditional average lies in that interval as well: for each
prefix length \(\nu\le n\), there are \(3^\nu\) equally weighted stream
assignments, and
\[
\Pr(K_0=\nu)\ge0,
\qquad
\sum_{\nu=0}^{n}\Pr(K_0=\nu)\le1,
\]
with the remaining probability assigned output zero. Thus
\[
\widehat\tau^{\mathrm{mix}}_{n,d,q}(\mathbf O_n)\in[-1,1],
\qquad
E_P\!\left[
(\widehat\tau^{\mathrm{mix}}_{n,d,q}-\Phi(P))^2
\right]\le4.
\]
If \(N\le1\), the auxiliary output and its conditional average are zero, so
\[
E_P\!\left[
(\widehat\tau^{\mathrm{mix}}_{n,d,q}-\Phi(P))^2
\right]=\Phi(P)^2\le1=\mathfrak U_{n,d,q}.
\]
If \(N>1\) and \(\mathfrak U_{n,d,q}=4\), the displayed bound by four proves
the required inequality. We henceforth consider
\(N>1\) and \(\mathfrak U_{n,d,q}\ne4\).
% lean: CausalSmith.Stat.MarRareqLogfrontier.unrestricted_auxiliary_risk_envelope

\item We first establish the transfer to independent Poisson streams.
Write
\[
\mathsf Q_{\mathrm{str}}
:=
\bigotimes_{h_{\mathrm{str}}\in\{0,1,2\}}
\left[
\sum_{\nu=0}^{\infty}
e^{-m}\frac{m^\nu}{\nu!}
\bigl(\mathcal L_P(O_1)\bigr)^{\otimes\nu}
\right],
\]
where each sum is a probability law on the disjoint union of finite
observed-record sequences. The stream labels \(0,1,2\) denote membership,
pilot, and estimation, respectively.

An uncapped prefix of length \(K_0\sim\operatorname{Poisson}(n/2)\), marked
independently with these three labels, has law \(\mathsf Q_{\mathrm{str}}\)
after separation into streams. Indeed, for any prescribed nonnegative
integer stream sizes \(\nu_0,\nu_1,\nu_2\), put
\(\nu_{\mathrm{tot}}:=\nu_0+\nu_1+\nu_2\). Their joint probability is
\[
e^{-n/2}\frac{(n/2)^{\nu_{\mathrm{tot}}}}{\nu_{\mathrm{tot}}!}
\frac{\nu_{\mathrm{tot}}!}{\nu_0!\nu_1!\nu_2!}
3^{-\nu_{\mathrm{tot}}}
=
\prod_{h_{\mathrm{str}}=0}^{2}
e^{-m}\frac{m^{\nu_{h_{\mathrm{str}}}}}
{\nu_{h_{\mathrm{str}}}!}.
\]
Conditional on these sizes, the records in the three streams have their
respective product observed-data laws. Equality on all count fibres proves
the asserted stream law.

Let
\[
\widetilde\tau_{\mathrm{Pois}}
:=
\operatorname{clip}_{[-1,1]}
\left(2\sum_{j=(a,x,s)\in\mathcal J_d}g(a)W_jG_j\right)
\]
be the uncapped output, using the statistics of
\cref{obj:synth_9,obj:synth_10,obj:synth_11,obj:synth_12,obj:synth_14,obj:synth_15,obj:synth_17}
on these finite streams. On the ratio branch, \(G_j=D_j\).
Here and below \(E_{\mathrm{Pois}}\) and
\(\Pr_{\mathrm{Pois}}\) denote expectation and probability under
\(\mathsf Q_{\mathrm{str}}\).

For the capped auxiliary output, conditional averaging gives
\[
\begin{aligned}
&E_\omega\!\left[
(\widetilde\tau(\mathbf O_n;\omega)-\Phi(P))^2
\mid\mathbf O_n\right]\\
&\quad=
(\widehat\tau^{\mathrm{mix}}_{n,d,q}(\mathbf O_n)-\Phi(P))^2
+
E_\omega\!\left[
(\widetilde\tau(\mathbf O_n;\omega)
-\widehat\tau^{\mathrm{mix}}_{n,d,q}(\mathbf O_n))^2
\mid\mathbf O_n\right].
\end{aligned}
\]
On each prefix-count fibre \(\nu\le n\), its stream law agrees with the
corresponding uncapped law. Summing these fibres, discarding the
nonnegative uncapped loss on the remaining fibres, and bounding the
overflow loss by four yields
\[
E_P\!\left[
(\widehat\tau^{\mathrm{mix}}_{n,d,q}-\Phi(P))^2
\right]
\le
E_{\mathrm{Pois}}\!\left[
(\widetilde\tau_{\mathrm{Pois}}-\Phi(P))^2
\right]+4\Pr(K_0>n).
\]
The Poisson probability series evaluates the exponential moment as
\[
E[2^{K_0}]
=e^{-n/2}\sum_{\nu=0}^{\infty}\frac{n^\nu}{\nu!}
=e^{n/2}.
\]
Since \(\{K_0>n\}\subseteq\{2^{K_0}\ge2^n\}\), exponential Markov gives
\[
\Pr(K_0>n)
\le\Pr(2^{K_0}\ge2^n)
\le2^{-n}E[2^{K_0}]
=e^{-n\log2+n/2}.
\]
Multiplying by \(e^{-n/2}\) cancels the intensity factor, so
\[
\Pr(K_0>n)e^{-n/2}
\le e^{-n\log2},
\qquad
\Pr(K_0>n)\le e^{-n\gamma_0}.
\]
Consequently,
\[
E_P\!\left[
(\widehat\tau^{\mathrm{mix}}_{n,d,q}-\Phi(P))^2
\right]
\le
E_{\mathrm{Pois}}\!\left[
(\widetilde\tau_{\mathrm{Pois}}-\Phi(P))^2
\right]+4e^{-n\gamma_0}.
\]
All these expectations are well defined; the capped sample space is finite
and the clipped outputs are bounded.
% lean: CausalSmith.Stat.MarRareqLogfrontier.deterministicRisk_mixedCountEstimator_le_streamRisk_add_tail

\item For every \(j=(a,x,s)\in\mathcal J_d\), define
\[
p_j:=p_{axs}(P),\qquad
\mu_j:=\mu_{axs}(P),\qquad
t_j:=mP(A=a,X=x,S=s,R=1).
\]
These definitions include null cells, with the mean convention of
\cref{obj:def:cell-functional}. Probability monotonicity and the
occupied-cell arrival condition give
\[
0\le\mu_j\le1,\qquad
t_j\ge N_0p_j,\qquad
\sum_jp_j=1,\qquad |\mathcal J_d|=4d.
\]
For \(p_j=0\), the membership, arrival, and arrived-one probabilities are
zero, and all the associated counts vanish almost surely.

Applying the same Poisson partition calculation to stream and cell
categories gives independent counts across distinct cells and streams,
with
\[
M_j\sim\operatorname{Poisson}(mp_j),
\qquad
C'_j,C_j\sim\operatorname{Poisson}(t_j).
\]
The estimation arrived-one and arrived-zero intensities are, respectively,
\(t_j\mu_j\) and \(t_j(1-\mu_j)\). Poisson moments therefore yield
\[
E_{\mathrm{Pois}}W_j=p_j,\qquad
E_{\mathrm{Pois}}W_j^2=p_j^2+\frac{p_j}{m}.
\]
Also, the membership statistic is independent of every statistic formed
from the pilot and estimation streams.
% lean: CausalSmith.Stat.MarRareqLogfrontier.integral_poissonMemberWeight_mul_selectedCellEstimate

\item Consider the polynomial branch
\(\ell\ge128\), \(d\ge\sqrt N\). Define the pilot-selection event by
\[
\mathsf L_j:=\{C'_j\le B/4\}.
\]
The coefficient and cell-statistic definitions in
\cref{obj:synth_13,obj:synth_14,obj:synth_15} give
\[
1-Q_k(t)=\sum_{v=1}^{k-1}a_vt^v,
\qquad
G_j=\mathbf1_{\mathsf L_j}H_j+
\mathbf1_{\mathsf L_j^c}D_j
\quad(t\ge0).
\]
To justify the expansion at zero and its degree, the numerator
\(1-T_k(1-2t/B)\) vanishes at zero and has degree at most \(k\).
The recurrence for \(T_k\), differentiated at one, gives
\[
T_0'(1)=0,\qquad T_1'(1)=1,\qquad
T_{h+1}'(1)=2+2T_h'(1)-T_{h-1}'(1)
\quad(h\ge1).
\]
Induction yields \(T_h'(1)=h^2\). Division of the numerator by \(t\)
therefore gives
\[
Q_k(0)=1,\qquad \deg(1-Q_k)\le k-1,
\]
which proves the displayed finite expansion.

For integers \(v\ge1\), write
\[
(C_j)_v:=\prod_{h=0}^{v-1}\max\{C_j-h,0\},
\qquad
\mathsf F_{j,v}:=
U_j\prod_{h=1}^{v-1}\max\{C_j-h,0\}.
\]
Empty products equal one. Counting ordered distinct tuples with the first
record an arrived one and the remaining records arrived gives
\[
E_{\mathrm{Pois}}\mathsf F_{j,v}
=t_j^v\mu_j.
\]
More explicitly, condition on the marked prefix size, count its ordered
distinct \(v\)-tuples, and then use the Poisson factorial moment
\(E(K_0)_v=(n/2)^v\). The one-record probabilities are
\(P(A=a,X=x,S=s,R=1,Y=1)/3\) for the distinguished record and
\(P(A=a,X=x,S=s,R=1)/3\) for the others. Their product gives the stated
identity. When the arrived-cell probability is zero, both sides vanish.
Thus
\[
E_{\mathrm{Pois}}H_j
=\mu_j\{1-Q_k(t_j)\}.
\]

For the ratio statistic, conditioning on the pattern of estimation
arrivals gives mean \(\mu_j\) when \(C_j>0\); its zero-count value is zero.
Consequently,
\[
E_{\mathrm{Pois}}D_j=\mu_j(1-e^{-t_j}),
\qquad 0\le D_j\le1.
\]
Pilot independence now gives the exact selected-statistic bias:
\[
E_{\mathrm{Pois}}G_j-\mu_j
=-\mu_j\left\{
\Pr_{\mathrm{Pois}}(\mathsf L_j)Q_k(t_j)
+\Pr_{\mathrm{Pois}}(\mathsf L_j^c)e^{-t_j}
\right\}.
\]
% lean: CausalSmith.Stat.MarRareqLogfrontier.markedPoissonSelectedCellEstimate_bias

\item We next derive the coefficient bound needed for the polynomial
moments. For a real polynomial \(F_{\mathrm{pol}}\), define
\[
\|F_{\mathrm{pol}}\|_{\mathrm{coef}}
:=\sum_{h\ge0}|[t_{\mathrm{pol}}^h]F_{\mathrm{pol}}|,
\]
where \(t_{\mathrm{pol}}\) is the polynomial indeterminate and the sum is
finite. For a second real polynomial and coefficient indices, take
\[
F_{\mathrm{in}}\in\mathbb R[t_{\mathrm{pol}}],
\qquad
h_{\mathrm{coef}},i_{\mathrm{coef}},j_{\mathrm{coef}}\in\mathbb N_0.
\]
All coefficient sums below are finite, with coefficients beyond the
polynomial degree equal to zero. The coefficientwise triangle inequality
 gives
\[
\begin{aligned}
\|F_{\mathrm{pol}}+F_{\mathrm{in}}\|_{\mathrm{coef}}
&=\sum_{h_{\mathrm{coef}}\ge0}
\left|[t_{\mathrm{pol}}^{h_{\mathrm{coef}}}]F_{\mathrm{pol}}
+[t_{\mathrm{pol}}^{h_{\mathrm{coef}}}]F_{\mathrm{in}}\right|\\
&\le\sum_{h_{\mathrm{coef}}\ge0}
\left|[t_{\mathrm{pol}}^{h_{\mathrm{coef}}}]F_{\mathrm{pol}}\right|
+\sum_{h_{\mathrm{coef}}\ge0}
\left|[t_{\mathrm{pol}}^{h_{\mathrm{coef}}}]F_{\mathrm{in}}\right|\\
&=\|F_{\mathrm{pol}}\|_{\mathrm{coef}}
+\|F_{\mathrm{in}}\|_{\mathrm{coef}}.
\end{aligned}
\]
Applying this inequality successively to the finite monomial expansion
of the product yields
\[
\begin{aligned}
\|F_{\mathrm{pol}}F_{\mathrm{in}}\|_{\mathrm{coef}}
&=\left\|
\sum_{i_{\mathrm{coef}}\ge0}\sum_{j_{\mathrm{coef}}\ge0}
\bigl([t_{\mathrm{pol}}^{i_{\mathrm{coef}}}]F_{\mathrm{pol}}\bigr)
\bigl([t_{\mathrm{pol}}^{j_{\mathrm{coef}}}]F_{\mathrm{in}}\bigr)
 t_{\mathrm{pol}}^{i_{\mathrm{coef}}+j_{\mathrm{coef}}}
\right\|_{\mathrm{coef}}\\
&\le\sum_{i_{\mathrm{coef}}\ge0}\sum_{j_{\mathrm{coef}}\ge0}
\left\|
\bigl([t_{\mathrm{pol}}^{i_{\mathrm{coef}}}]F_{\mathrm{pol}}\bigr)
\bigl([t_{\mathrm{pol}}^{j_{\mathrm{coef}}}]F_{\mathrm{in}}\bigr)
 t_{\mathrm{pol}}^{i_{\mathrm{coef}}+j_{\mathrm{coef}}}
\right\|_{\mathrm{coef}}\\
&=\sum_{i_{\mathrm{coef}}\ge0}\sum_{j_{\mathrm{coef}}\ge0}
\left|[t_{\mathrm{pol}}^{i_{\mathrm{coef}}}]F_{\mathrm{pol}}\right|
\left|[t_{\mathrm{pol}}^{j_{\mathrm{coef}}}]F_{\mathrm{in}}\right|\\
&=\left(\sum_{i_{\mathrm{coef}}\ge0}
\left|[t_{\mathrm{pol}}^{i_{\mathrm{coef}}}]F_{\mathrm{pol}}\right|\right)
\left(\sum_{j_{\mathrm{coef}}\ge0}
\left|[t_{\mathrm{pol}}^{j_{\mathrm{coef}}}]F_{\mathrm{in}}\right|\right)\\
&=\|F_{\mathrm{pol}}\|_{\mathrm{coef}}
\|F_{\mathrm{in}}\|_{\mathrm{coef}}.
\end{aligned}
\]
Here the norm of each monomial is the absolute value of its coefficient,
by the defining coefficient sum. For every real polynomial \(F_{\mathrm{in}}\), submultiplicativity gives by induction
\[
\|F_{\mathrm{in}}^h\|_{\mathrm{coef}}
\le\|F_{\mathrm{in}}\|_{\mathrm{coef}}^h
\qquad(h\in\mathbb N_0).
\]
For \(F_{\mathrm{pol}}\ne0\), expand the composition as the finite sum
\[
F_{\mathrm{pol}}\circ F_{\mathrm{in}}
=
\sum_{h=0}^{\deg F_{\mathrm{pol}}}
\bigl([t_{\mathrm{pol}}^h]F_{\mathrm{pol}}\bigr)F_{\mathrm{in}}^h.
\]
For the composition bound, assume
\(\|F_{\mathrm{in}}\|_{\mathrm{coef}}\ge1\).
Subadditivity and the power bound then give
\[
\begin{aligned}
\|F_{\mathrm{pol}}\circ F_{\mathrm{in}}\|_{\mathrm{coef}}
&\le
\sum_{h=0}^{\deg F_{\mathrm{pol}}}
|[t_{\mathrm{pol}}^h]F_{\mathrm{pol}}|
\|F_{\mathrm{in}}^h\|_{\mathrm{coef}}\\
&\le
\sum_{h=0}^{\deg F_{\mathrm{pol}}}
|[t_{\mathrm{pol}}^h]F_{\mathrm{pol}}|
\|F_{\mathrm{in}}\|_{\mathrm{coef}}^h\\
&\le
\sum_{h=0}^{\deg F_{\mathrm{pol}}}
|[t_{\mathrm{pol}}^h]F_{\mathrm{pol}}|
\|F_{\mathrm{in}}\|_{\mathrm{coef}}^{\deg F_{\mathrm{pol}}}\\
&=
\|F_{\mathrm{pol}}\|_{\mathrm{coef}}
\|F_{\mathrm{in}}\|_{\mathrm{coef}}^{\deg F_{\mathrm{pol}}}.
\end{aligned}
\]
When \(F_{\mathrm{pol}}=0\), the composition has coefficient norm zero.
The Chebyshev recurrence implies
\[
\|T_{h+2}\|_{\mathrm{coef}}
\le2\|T_{h+1}\|_{\mathrm{coef}}+\|T_h\|_{\mathrm{coef}},
\qquad
\|T_0\|_{\mathrm{coef}}=\|T_1\|_{\mathrm{coef}}=1.
\]
Since \((1+\sqrt2)^2=2(1+\sqrt2)+1\), induction gives
\[
\|T_h\|_{\mathrm{coef}}\le(1+\sqrt2)^h
\quad(h\in\mathbb N).
\]
Define
\[
\mathcal P_k(t_{\mathrm{pol}})
:=1-T_k(1-2t_{\mathrm{pol}}).
\]
Then
\[
\|\mathcal P_k\|_{\mathrm{coef}}
\le1+(1+\sqrt2)^k3^k\le1+8^k,
\]
using \(\|1-2t_{\mathrm{pol}}\|_{\mathrm{coef}}=3\) and
\(\sqrt2\le3/2\). The scaled quotient satisfies
\[
Q_k(Bt_{\mathrm{pol}})
=\frac1{2k^2}\,
\frac{\mathcal P_k(t_{\mathrm{pol}})}{t_{\mathrm{pol}}},
\]
where division is exact because \(\mathcal P_k(0)=0\).
Removing the zero constant coefficient cannot increase the coefficient
norm. Hence, using \(k\ge2\) and \(8^k\ge64\),
\[
\begin{aligned}
\|1-Q_k(Bt_{\mathrm{pol}})\|_{\mathrm{coef}}
&\le1+\frac{1+8^k}{2k^2}\\
&\le1+\frac{1+8^k}{8}
\le8^k.
\end{aligned}
\]
In particular,
\[
\sum_{v=1}^{k-1}|a_vB^v|\le8^k.
\]
% lean: CausalSmith.Stat.MarRareqLogfrontier.needle_scaled_coefficient_oneNorm_le

\item For \(t\ge0\) and integers \(v,w\ge1\), define the factorial overlap
\[
\mathcal A_{v,w}(t)
:=
\sum_{h=0}^{\min\{v,w\}}
\binom vh\binom wh h!\,t^{v+w-h}.
\]
Counting the \(h\) shared entries of two ordered distinct tuples gives
\[
(C_j)_v(C_j)_w
=
\sum_{h=0}^{\min\{v,w\}}
\binom vh\binom wh h!(C_j)_{v+w-h}.
\]
Taking Poisson factorial moments yields
\[
E_{\mathrm{Pois}}[(C_j)_v(C_j)_w]
=\mathcal A_{v,w}(t_j).
\]
Since \(0\le\mathsf F_{j,v}\le(C_j)_v\), expansion of \(H_j^2\) gives
\[
E_{\mathrm{Pois}}H_j^2
\le
\sum_{v=1}^{k-1}\sum_{w=1}^{k-1}
|a_va_w|\mathcal A_{v,w}(t_j).
\]

For \(0\le t\le B\) and \(v,w\le k\), the bound
\[
\binom vh\binom wh h!\le\frac{(vw)^h}{h!}
\]
implies
\[
\begin{aligned}
\mathcal A_{v,w}(t)
&\le B^{v+w}
\sum_{h=0}^{\min\{v,w\}}
\frac{(k^2/B)^h}{h!}\\
&\le B^{v+w}e^{k^2/B}.
\end{aligned}
\]
Therefore, when \(t_j\le B\),
\[
E_{\mathrm{Pois}}H_j^2
\le e^{k^2/B}
\left(\sum_{v=1}^{k-1}|a_vB^v|\right)^2
\le e^{k^2/B}64^k.
\]
The tuning bounds give
\[
\frac{k^2}{B}\le\frac{\ell}{64},
\qquad 64^k=2^{6k}\le e^{6k},
\]
so
\[
E_{\mathrm{Pois}}H_j^2
\le e^{k^2/B+6k}\le e^{\ell/8}
\quad(t_j\le B).
\]
% lean: CausalSmith.Stat.MarRareqLogfrontier.integral_sq_markedPoissonFactorialBranch_le_exp_of_light

\item A second overlap estimate controls arbitrarily large intensities.
Using \(\binom whh!\le w^h\) and the binomial theorem,
\[
\begin{aligned}
\mathcal A_{v,w}(t)
&\le
\sum_{h=0}^{v}\binom vh w^h t^{v-h}t^w\\
&=(t+w)^vt^w
\le(t+v+w)^{v+w}.
\end{aligned}
\]
For \(v,w\le k\), split according as \((t+2k)/B\le1\) or exceeds one.
In both cases,
\[
\left(\frac{t+2k}{B}\right)^{v+w}
\le
\max\left\{1,\left(\frac{t+2k}{B}\right)^{2k}\right\}.
\]
The coefficient bound thus gives
\[
E_{\mathrm{Pois}}H_j^2
\le64^k
\max\left\{1,
\left(\frac{t_j+2k}{B}\right)^{2k}\right\}.
\]
% lean: CausalSmith.Stat.MarRareqLogfrontier.integral_sq_markedPoissonFactorialBranch_le_full_growth

\item Suppose \(t_j>B\). Define the integer pilot cutoff by
\[
\nu_{\mathrm{cut}}:=\lfloor B/4\rfloor.
\]
Then \(\nu_{\mathrm{cut}}<t_j/4\). The Poisson lower-tail estimate used here
follows from the centered exponential moment. To see its relevant form,
put
\[
a_{\mathrm{tail}}:=\frac1{\sqrt2}.
\]
The elementary exponential inequality
\[
e^{-a_{\mathrm{tail}}}
\le1-a_{\mathrm{tail}}+\frac{a_{\mathrm{tail}}^2}{2}
\]
follows by multiplying
\(e^{a_{\mathrm{tail}}}\ge
1+a_{\mathrm{tail}}+a_{\mathrm{tail}}^2/2\) by the nonnegative quadratic
on the right and using
\[
\left(1+a_{\mathrm{tail}}+\frac{a_{\mathrm{tail}}^2}{2}\right)
\left(1-a_{\mathrm{tail}}+\frac{a_{\mathrm{tail}}^2}{2}\right)
=1+\frac{a_{\mathrm{tail}}^4}{4}\ge1.
\]
For \(C'_j\sim\operatorname{Poisson}(t_j)\), exponential Markov therefore
gives
\[
\begin{aligned}
\Pr_{\mathrm{Pois}}\left(t_j-C'_j>\frac{t_j}{\sqrt2}\right)
&\le
\exp\left\{-\frac{t_j}{2}
+t_j(e^{-a_{\mathrm{tail}}}-1+a_{\mathrm{tail}})\right\}\\
&\le e^{-t_j/4}.
\end{aligned}
\]
On \(\{C'_j\le\nu_{\mathrm{cut}}\}\),
\(t_j-C'_j>3t_j/4>t_j/\sqrt2\). Consequently,
\[
\Pr_{\mathrm{Pois}}(\mathsf L_j)\le e^{-t_j/4}.
\]

Since \(t_j>B\), the maximum in the preceding moment bound selects its
second argument. The inequalities \(64^k\le e^{6k}\) and
\(t_{\mathrm{grow}}\le e^{t_{\mathrm{grow}}}\) for
\(t_{\mathrm{grow}}:=(t_j+2k)/B\) give
\[
\Pr_{\mathrm{Pois}}(\mathsf L_j)E_{\mathrm{Pois}}H_j^2
\le
\exp\left\{-\frac{t_j}{4}
+6k+\frac{2k(t_j+2k)}{B}\right\}.
\]
Here
\[
k\le\frac{t_j}{16384},
\qquad \frac{2k}{B}\le\frac1{8192},
\]
and hence
\[
6k+\frac{2k(t_j+2k)}B
\le
t_j\left(\frac6{16384}+\frac1{8192}
+\frac1{67108864}\right)
\le\frac{t_j}{8}.
\]
Thus
\[
\Pr_{\mathrm{Pois}}(\mathsf L_j)E_{\mathrm{Pois}}H_j^2
\le e^{-t_j/8}\le e^{-32\ell}.
\]
Also \(\Pr_{\mathrm{Pois}}(\mathsf L_j)\le e^{-32\ell}\).
% lean: CausalSmith.Stat.MarRareqLogfrontier.markedPoisson_pilot_light_mul_factorial_secondMoment_le_intensity

\item We now bound the two terms in the selected-statistic bias.
For \(0\le t\le B\), the polynomial quotient in
\cref{obj:synth_18} satisfies
\[
tQ_k(t)=\frac{B}{2k^2}\{1-T_k(1-2t/B)\}.
\]
Because \(|1-2t/B|\le1\), the cosine characterization of \(T_k\) gives
\(|T_k(1-2t/B)|\le1\), and therefore
\[
t|Q_k(t)|\le\frac B{k^2}.
\]
For \(t_j\le B\), use \(N_0p_j\mu_j\le t_j\) to obtain
\[
p_j\mu_j
\Pr_{\mathrm{Pois}}(\mathsf L_j)|Q_k(t_j)|
\le\frac{B}{N_0k^2}.
\]
This includes \(t_j=0\), because then \(p_j=0\).

For \(t_j>B\), Cauchy--Schwarz and the heavy moment estimate give
\[
\begin{aligned}
\left\{
\Pr_{\mathrm{Pois}}(\mathsf L_j)
|E_{\mathrm{Pois}}H_j|
\right\}^2
&\le
\Pr_{\mathrm{Pois}}(\mathsf L_j)
E_{\mathrm{Pois}}H_j^2\\
&\le e^{-t_j/8}.
\end{aligned}
\]
Using
\(\mu_jQ_k(t_j)=\mu_j-E_{\mathrm{Pois}}H_j\), it follows that
\[
\begin{aligned}
\mu_j\Pr_{\mathrm{Pois}}(\mathsf L_j)|Q_k(t_j)|
&\le
\Pr_{\mathrm{Pois}}(\mathsf L_j)
+\Pr_{\mathrm{Pois}}(\mathsf L_j)|E_{\mathrm{Pois}}H_j|\\
&\le2e^{-t_j/16}\le2e^{-16\ell}.
\end{aligned}
\]
Splitting the cell sum at \(t_j=B\), using \(|g(a)|=1\) from
\cref{obj:synth_19}, and then using \(4d\) cells and \(\sum_jp_j=1\), yields
\[
\left|
2\sum_jg(a)p_j\mu_j
\Pr_{\mathrm{Pois}}(\mathsf L_j)Q_k(t_j)
\right|
\le\frac{8dB}{N_0k^2}+4e^{-16\ell}.
\]
% lean: CausalSmith.Stat.MarRareqLogfrontier.abs_sum_pilot_needle_bias_le

\item The complementary-pilot term uses an upper-tail bound with the same
integer cutoff. The Poisson probability series gives
\[
E_{\mathrm{Pois}}[2^{C'_j}]
=e^{-t_j}\sum_{\nu=0}^{\infty}\frac{(2t_j)^\nu}{\nu!}
=e^{t_j}.
\]
The threshold inclusion and exponential Markov inequality imply
\[
\begin{aligned}
\Pr_{\mathrm{Pois}}(C'_j>\nu_{\mathrm{cut}})
&\le\Pr_{\mathrm{Pois}}(2^{C'_j}\ge2^{\nu_{\mathrm{cut}}})\\
&\le2^{-\nu_{\mathrm{cut}}}E_{\mathrm{Pois}}[2^{C'_j}]
=e^{-\nu_{\mathrm{cut}}\log2+t_j}.
\end{aligned}
\]
Multiplying by \(e^{-t_j}\) cancels the intensity factor and yields
\[
\Pr_{\mathrm{Pois}}(C'_j>\nu_{\mathrm{cut}})e^{-t_j}
\le e^{-\nu_{\mathrm{cut}}\log2}.
\]
Since
\[
64\ell<\nu_{\mathrm{cut}}+1,\qquad
\ell\ge128,\qquad \log2>\frac58,
\]
we have \(\nu_{\mathrm{cut}}\log2\ge16\ell\). Thus
\[
\Pr_{\mathrm{Pois}}(\mathsf L_j^c)e^{-t_j}\le e^{-16\ell},
\]
including \(t_j=0\), when the upper-tail probability is zero.
Consequently,
\[
\left|
2\sum_jg(a)p_j\mu_j
\Pr_{\mathrm{Pois}}(\mathsf L_j^c)e^{-t_j}
\right|
\le2e^{-16\ell}.
\]

Define the ideal raw contrast by
\[
\tau_{\mathrm{raw}}
:=2\sum_jg(a)W_jG_j.
\]
Membership independence gives
\(E_{\mathrm{Pois}}(W_jG_j)=p_jE_{\mathrm{Pois}}G_j\).
Combining the exact bias identity with the last two aggregate bounds,
\[
|E_{\mathrm{Pois}}\tau_{\mathrm{raw}}-\Phi(P)|
\le
\frac{8dB}{N_0k^2}+6e^{-16\ell}
=b_{n,d,q}.
\]
% lean: CausalSmith.Stat.MarRareqLogfrontier.abs_integral_raw_poissonSelected_estimator_bias_le

\item We next establish the ratio variance contribution.
For a fixed marked prefix length \(\nu\), define the estimation-arrival
index set by
\[
\mathsf I_j^{(\nu)}
:=
\{i\in\{1,\ldots,\nu\}:
i\text{ has estimation label and is an arrival in cell }j\}.
\]
On a prescribed pattern \(\mathsf I_j^{(\nu)}=\mathsf I_{\mathrm{arr}}\) with
\(\nu_{\mathrm{arr}}:=|\mathsf I_{\mathrm{arr}}|>0\), the restricted product law makes the
arrived-outcome indicators independent with mean \(\mu_j\). Their centered
first moments vanish, and their centered squares are at most one.
Expanding the square therefore gives, for patterns of positive
probability,
\[
E\!\left[
(U_j-\nu_{\mathrm{arr}}\mu_j)^2
\mid\mathsf I_j^{(\nu)}=\mathsf I_{\mathrm{arr}}
\right]\le\nu_{\mathrm{arr}},
\]
and hence
\[
E\!\left[
(D_j-\mu_j)^2
\mid\mathsf I_j^{(\nu)}=\mathsf I_{\mathrm{arr}}
\right]\le\nu_{\mathrm{arr}}^{-1}.
\]
Zero-probability patterns contribute zero. On the zero-count pattern,
\((D_j-\mu_j)^2=\mu_j^2\le1\). Summing first over arrival patterns and then
over the Poisson prefix size gives
\[
E_{\mathrm{Pois}}(D_j-\mu_j)^2
\le
E_{\mathrm{Pois}}\!\left[
\mathbf1_{\{C_j=0\}}
+\frac{\mathbf1_{\{C_j>0\}}}{C_j}
\right],
\]
where the quotient is zero at \(C_j=0\).

If \(t_j\le1\), the integrand on the right is at most one, so its expectation
is at most \(1\le4/(1+t_j)\). If \(t_j>1\), its value is at most
\(2/(C_j+1)\). For every nonnegative integer \(h_{\mathrm{count}}\), the Poisson probabilities satisfy
\[
\Pr_{\mathrm{Pois}}(C_j=h_{\mathrm{count}})
=e^{-t_j}\frac{t_j^{h_{\mathrm{count}}}}{h_{\mathrm{count}}!},
\qquad
\frac{t_j}{h_{\mathrm{count}}+1}
\Pr_{\mathrm{Pois}}(C_j=h_{\mathrm{count}})
=
\Pr_{\mathrm{Pois}}(C_j=h_{\mathrm{count}}+1).
\]
The probability series sums to one, and the reciprocal series is bounded termwise by it. Summing the recurrence and shifting the index therefore gives
\[
\begin{aligned}
t_jE_{\mathrm{Pois}}\frac1{C_j+1}
&=
\sum_{h_{\mathrm{count}}=0}^{\infty}
\frac{t_j}{h_{\mathrm{count}}+1}
\Pr_{\mathrm{Pois}}(C_j=h_{\mathrm{count}})\\
&=
\sum_{h_{\mathrm{count}}=0}^{\infty}
\Pr_{\mathrm{Pois}}(C_j=h_{\mathrm{count}}+1)\\
&=\Pr_{\mathrm{Pois}}(C_j>0)
=1-e^{-t_j}.
\end{aligned}
\]
Dividing by \(t_j>1\) yields
\[
E_{\mathrm{Pois}}\frac1{C_j+1}
=\frac{1-e^{-t_j}}{t_j}.
\]
Therefore
\[
E_{\mathrm{Pois}}(D_j-\mu_j)^2
\le\frac2{t_j}\le\frac4{1+t_j}
\quad(t_j>1).
\]
When the arrived-cell probability is zero, \(C_j=U_j=\mu_j=0\)
almost surely, so the centered second moment is zero. In all cases,
\[
\operatorname{Var}_{\mathrm{Pois}}(D_j)
=
\operatorname{Var}_{\mathrm{Pois}}(D_j-\mu_j)
\le E_{\mathrm{Pois}}(D_j-\mu_j)^2
\le\frac4{1+t_j}.
\]

Expanding the variance of the independent product \(W_jD_j\) gives
\[
\operatorname{Var}_{\mathrm{Pois}}(W_jD_j)
=p_j^2\operatorname{Var}_{\mathrm{Pois}}(D_j)
+\frac{p_j}{m}E_{\mathrm{Pois}}D_j^2.
\]
Using \(D_j^2\le1\) and \(N_0p_j\le t_j\),
\[
\frac{p_j^2}{1+t_j}\le\frac{p_j}{N_0},
\]
including \(p_j=0\). Summing over cells proves
\[
\sum_j\operatorname{Var}_{\mathrm{Pois}}(W_jD_j)
\le\frac4{N_0}+\frac1m.
\]
% lean: CausalSmith.Stat.MarRareqLogfrontier.sum_variance_poissonMemberWeight_mul_ratioBranch_le

\item On the polynomial branch, define the weighted correction by
\[
\Delta_j:=W_j(G_j-D_j)
=W_j\mathbf1_{\mathsf L_j}(H_j-D_j).
\]
Pilot and membership independence give
\[
E_{\mathrm{Pois}}\Delta_j^2
=
\left(p_j^2+\frac{p_j}{m}\right)
\Pr_{\mathrm{Pois}}(\mathsf L_j)
E_{\mathrm{Pois}}(H_j-D_j)^2.
\]
Since \((H_j-D_j)^2\le2(H_j^2+D_j^2)\), for \(t_j\le B\),
\[
E_{\mathrm{Pois}}\Delta_j^2
\le2\left(p_j^2+\frac{p_j}{m}\right)(e^{\ell/8}+1).
\]
Here \(p_j\le B/N_0\), so
\[
p_j^2+\frac{p_j}{m}
\le\frac{B^2}{N_0^2}+\frac{B}{mN_0}.
\]
Summing over at most \(4d\) cells,
\[
8\sum_{j:t_j\le B}E_{\mathrm{Pois}}\Delta_j^2
\le
64d(e^{\ell/8}+1)
\left(\frac{B^2}{N_0^2}+\frac{B}{mN_0}\right).
\]
For \(t_j>B\), the two heavy-pilot estimates instead yield
\[
\Pr_{\mathrm{Pois}}(\mathsf L_j)
E_{\mathrm{Pois}}(H_j-D_j)^2
\le4e^{-32\ell},
\]
and therefore
\[
8\sum_{j:t_j>B}E_{\mathrm{Pois}}\Delta_j^2
\le32e^{-32\ell}\left(1+\frac1m\right).
\]
In the last step we used
\[
\sum_j\left(p_j^2+\frac{p_j}{m}\right)
\le1+\frac1m,
\]
because \(0\le p_j\le1\) and \(\sum_jp_j=1\).
% lean: CausalSmith.Stat.MarRareqLogfrontier.sum_integral_sq_memberWeight_mul_selectedCorrection_le

\item Distinct cells are independent in the uncapped marked experiment.
Since \(g(a)^2=1\),
\[
\operatorname{Var}_{\mathrm{Pois}}(\tau_{\mathrm{raw}})
=4\sum_j\operatorname{Var}_{\mathrm{Pois}}(W_jG_j).
\]
For each cell, \(W_jG_j=W_jD_j+\Delta_j\). Nonnegativity of
\(\operatorname{Var}_{\mathrm{Pois}}(W_jD_j-\Delta_j)\) implies
\[
\operatorname{Var}_{\mathrm{Pois}}(W_jG_j)
\le
2\operatorname{Var}_{\mathrm{Pois}}(W_jD_j)
+2\operatorname{Var}_{\mathrm{Pois}}(\Delta_j)
\le
2\operatorname{Var}_{\mathrm{Pois}}(W_jD_j)
+2E_{\mathrm{Pois}}\Delta_j^2.
\]
The preceding aggregate estimates thus give
\[
\operatorname{Var}_{\mathrm{Pois}}(\tau_{\mathrm{raw}})
\le v_{n,d,q}.
\]
The raw contrast has finite second moment by the finite polynomial moment
bounds already established, and its squared error decomposes as
\[
E_{\mathrm{Pois}}(\tau_{\mathrm{raw}}-\Phi(P))^2
=
\operatorname{Var}_{\mathrm{Pois}}(\tau_{\mathrm{raw}})
+(E_{\mathrm{Pois}}\tau_{\mathrm{raw}}-\Phi(P))^2
\le v_{n,d,q}+b_{n,d,q}^2.
\]
Clipping decreases squared distance to every point of \([-1,1]\): below
\(-1\) it replaces the value by \(-1\), above \(1\) by \(1\), and inside the
interval it preserves the value. Since \(\Phi(P)\in[-1,1]\),
\[
E_{\mathrm{Pois}}(\widetilde\tau_{\mathrm{Pois}}-\Phi(P))^2
\le b_{n,d,q}^2+v_{n,d,q}.
\]
Combining this with the prefix transfer and the bound by four proves
\[
E_P\!\left[
(\widehat\tau^{\mathrm{mix}}_{n,d,q}-\Phi(P))^2
\right]
\le
\min\{4,b_{n,d,q}^2+v_{n,d,q}+4e^{-n\gamma_0}\}
=\mathfrak U_{n,d,q}.
\]
% lean: CausalSmith.Stat.MarRareqLogfrontier.unrestricted_auxiliary_risk_envelope_active

\item Finally, consider the ratio branch. Its ideal raw contrast is
\[
\tau_{\mathrm{raw}}^{\mathrm{ratio}}
:=2\sum_jg(a)W_jD_j.
\]
The ratio mean and membership independence give
\[
E_{\mathrm{Pois}}\tau_{\mathrm{raw}}^{\mathrm{ratio}}-\Phi(P)
=-2\sum_jg(a)p_j\mu_je^{-t_j}.
\]
Because \(t_j\ge N_0p_j\), \(0\le\mu_j\le1\), and
\(t_{\mathrm{bias}}e^{-t_{\mathrm{bias}}}\le e^{-1}\) for every real
\(t_{\mathrm{bias}}\ge0\),
\[
p_j\mu_je^{-t_j}
\le p_je^{-N_0p_j}
\le\frac1{eN_0}.
\]
The exponential inequality follows from
\(e^{t_{\mathrm{bias}}-1}\ge t_{\mathrm{bias}}\).
Thus
\[
|E_{\mathrm{Pois}}\tau_{\mathrm{raw}}^{\mathrm{ratio}}-\Phi(P)|
\le\frac{8d}{eN_0}.
\]
Independence across cells and the ratio variance aggregate give
\[
\operatorname{Var}_{\mathrm{Pois}}
(\tau_{\mathrm{raw}}^{\mathrm{ratio}})
=
4\sum_j\operatorname{Var}_{\mathrm{Pois}}(W_jD_j)
\le4\left(\frac4{N_0}+\frac1m\right).
\]
The squared-error decomposition, clipping, prefix transfer, and bound by
four therefore imply
\[
\begin{aligned}
&E_P\!\left[
(\widehat\tau^{\mathrm{mix}}_{n,d,q}-\Phi(P))^2
\right]\\
&\quad\le
\min\left\{4,
\left(\frac{8d}{eN_0}\right)^2
+4\left(\frac4{N_0}+\frac1m\right)
+4e^{-n\gamma_0}\right\}
=\mathfrak U_{n,d,q}.
\end{aligned}
\]
Together with the zero and capped-envelope cases, this establishes the
risk bound for every admissible \(P\). Combining it with the deterministic
comparison proved above gives
\[
E_P\!\left[
\left(\widehat\tau^{\mathrm{mix}}_{n,d,q}(\mathbf O_n)-\Phi(P)\right)^2
\right]
\le\mathfrak U_{n,d,q}
\le10^{30}r_{n,d,q}.
\]
The expectation is under the canonical product observed-sample law
specified in the statement, and \(10^{30}\) is positive, finite, and
independent of \(P,n,d,q\).
% lean: CausalSmith.Stat.MarRareqLogfrontier.unrestricted_uniform_mixed_count_certificate
\end{enumerate}
\end{proof}
